\begin{statementbox}
	\textbf{\textcolor{CStatement}{条件}}
	\begin{description}[style=nextline, leftmargin=2.8em, labelsep=0.5em, itemsep=0.35em]
		\item[\textbf{\textcolor{CStatement}{条件 1（曲线与点）：}}]
			已知非退化圆锥曲线的标准方程（椭圆 $\dfrac{x^2}{a^2}+\dfrac{y^2}{b^2}=1$、双曲线 $\dfrac{x^2}{a^2}-\dfrac{y^2}{b^2}=1$、抛物线 $y^2=2px$）和平面内一点 $P(x_0,y_0)$。
	\end{description}
	\textbf{\textcolor{CStatement}{结论}}
	\begin{description}[style=nextline, leftmargin=2.8em, labelsep=0.5em, itemsep=0.45em]
		\item[\textbf{\textcolor{CStatement}{结论一（极线方程）：}}]
			\textbf{\textcolor{CStatement}{直观描述：}}
			将曲线方程中的平方项替换为乘积项、一次项替换为平均值，即得极线方程。当点 $P$ 在曲线上时，该极线为切线；当点 $P$ 在曲线外时，该极线为切点弦。具体形式如下：\textbf{椭圆：}
			\[
				\frac{x_0x}{a^2}+\frac{y_0y}{b^2}=1
			\]
			\textbf{双曲线：}
			\[
				\frac{x_0x}{a^2}-\frac{y_0y}{b^2}=1
			\]
			\textbf{抛物线：}
			\[
				y_0y=p(x+x_0)
			\]
		\item[\textbf{\textcolor{CStatement}{推广结论（配极原理）：}}]
			\textbf{\textcolor{CStatement}{直观描述：}}
			极点与极线的对称关系：若点 $Q$ 在点 $P$ 的极线上，则点 $P$ 也在点 $Q$ 的极线上，反之亦然。
			\[
				Q\in l_P \iff P\in l_Q .
			\]
	\end{description}

\end{statementbox}
